%%%%%%%%%%%%%%%%%%%%%%%%%%%%%%%%%%%%%%%%%%%%%%%%%%%%%%%%%%%%%%%%%%%%%%%%%%%%%%%
%
% Präambel
%
%%%%%%%%%%%%%%%%%%%%%%%%%%%%%%%%%%%%%%%%%%%%%%%%%%%%%%%%%%%%%%%%%%%%%%%%%%%%%%%
%\RequirePackage{pdf14}
\newcommand{\hh}{20mm}                    %% Höhe der Kopfzeile
\documentclass[
draft=false,
paper=A4,
pagesize,
headheight=\hh,
DIV=15,
BCOR=10mm,
twoside=false,
headinclude=true,
footinclude=false,
fontsize=12pt,
numbers=autoenddot,
toc=bibliography
]{scrreprt}
\usepackage[english]{babel}
\usepackage[table]{xcolor}
\usepackage{pdflscape}
\usepackage{rotating}
\usepackage{chngcntr}
\counterwithout{table}{chapter}
\usepackage{caption}

\usepackage{array}
\newcolumntype{L}[1]{>{\raggedright\let\newline\\\arraybackslash\hspace{0pt}}m{#1}}
\newcolumntype{C}[1]{>{\centering\let\newline\\\arraybackslash\hspace{0pt}}m{#1}}
\newcolumntype{R}[1]{>{\raggedleft\let\newline\\\arraybackslash\hspace{0pt}}m{#1}}

% % % INCREASING THE NUMBER OF POSSIBLE FLOATS... % %
\usepackage{etex}
\reserveinserts{36}
\usepackage{morefloats}
% % % INCREASING THE NUMBER OF POSSIBLE FLOATS... % %

\usepackage[nomessages]{fp}
\usepackage{lipsum}
\usepackage{tocloft}
\usepackage{caption}
\usepackage{subcaption}

% Schriftarten %%%%%%%%%%%%%%%%%%%%%%%%%%%%%%%%%%%%%%%%%%%%%%%%%%%%%%%%%%%%%%%%

\usepackage{mathptmx}                          %% serif      = Times
\usepackage[scaled=.91]{helvet}                %% sans serif = Helvetica
\usepackage{courier}                           %% typewriter = Courier
\addtokomafont{disposition}{\rmfamily}         %% Änderung Überschriften-Font
\addtokomafont{descriptionlabel}{\rmfamily}    %% Änderung Description-Font


%%%%%%%%%%%%%%%%%%%%%%%%%%%%%%%%%%%%%%%%%%%%%%%%%%%%%%%%%%%%%%%%%%%%%%%%%%%%%%%


\usepackage[T1]{fontenc}
\usepackage[ansinew]{inputenc}
\usepackage[babel]{microtype}
\usepackage{graphicx}
\usepackage{arydshln}
\usepackage[space]{grffile}
%% \usepackage{epic, eepic}
\usepackage{rotating}
\usepackage{lastpage}
\usepackage[numbers, sort]{natbib}
\usepackage{csquotes}
\usepackage[section,above,below]{placeins}
\usepackage{amssymb,amsmath}
\usepackage{enumitem}
\usepackage{pdfpages}

\usepackage{caption}
\usepackage{subcaption}

\usepackage{tikz}

% Randnotizen %%%%%%%%%%%%%%%%%%%%%%%%%%%%%%%%%%%%%%%%%%%%%%%%%%%%%%%%%%%%%%%%%

\usepackage{marginnote}
\usepackage{ragged2e}
\usepackage{setspace}
\usepackage{ifthen}
\renewcommand*{\raggedleftmarginnote}{\RaggedLeft}
\renewcommand*{\raggedrightmarginnote}{\RaggedRight}
\reversemarginpar
\setlength{\marginparwidth}{1.75\marginparwidth}

\newcommand{\mnote}[2][\empty]{%
  \marginnote{%
    \begin{spacing}{0.9}
      \ifthenelse{\equal{#1}{\empty}} {} {\color{#1}}
      \vspace{-\baselineskip}\vspace{-\parskip}\scriptsize\hspace{0pt}#2%
    \end{spacing}
  }
}

\usepackage{siunitx}
\usepackage{float}
\usepackage{arydshln}

%%%%%%%%%%%%%%%%%%%%%%%%%%% SONSTIGE AENDERUNGEN %%%%%%%%%%%%%%%%%%%%%%%%%%%%%%%
\usepackage{datatool}
% Define a convenient command to add a line
% to the database
\newcommand*{\addacronym}[2]{%
	\DTLnewrow{acronyms}%
	\DTLnewdbentry{acronyms}{Acronym}{#1}%
	\DTLnewdbentry{acronyms}{Description}{#2}%
}
%% Pakage and command for inserting a blank page
\usepackage{afterpage}
\newcommand\blankpage{%
	\null
	\thispagestyle{empty}%
	\addtocounter{page}{-1}%
	\newpage}
% Header and Footer (Style "TUB") %%%%%%%%%%%%%%%%%%%%%%%%%%%%%%%%%%%%%%%%%%%%

\usepackage[headsepline]{scrlayer-scrpage}


\setkomafont{pageheadfoot}{\normalfont\mdseries\upshape}
\setkomafont{pagenumber}{\normalfont\mdseries\upshape}

\newlength{\A}
\newlength{\B}
\newlength{\C}
\newlength{\D}
\newlength{\E}
\newlength{\F}
\setlength{\C}{0.36\textwidth}
\setlength{\A}{\hh}
\setlength{\A}{1.10\A}
\setlength{\B}{0.31\textwidth}
\addtolength{\B}{-\A} 
\setlength{\D}{0.11\textwidth}
\setlength{\E}{0.22\textwidth}
\setlength{\F}{\hh}
\setlength{\F}{0.75\F}


\newcommand{\thisauthor}{Name}
\newpagestyle{TUB}
{%% Kopfzeile
  (\textwidth, 0pt)
  {} {}
  {%
    \begin{minipage}[c][\hh]{\A}
      \includegraphics[height=\F]{tuBerlinLogoSchwarz}
    \end{minipage}%
    \begin{minipage}[c][\hh]{\B}
      \flushleft\scriptsize
      \textsf{Institut f\"ur\\Luft- und Raumfahrt\\[4pt]
        Fachgebiet\\Raumfahrttechnik}
    \end{minipage}%
    \begin{minipage}[c][\hh]{\C}
      \centering\footnotesize\textbf\st
    \end{minipage}%
    \begin{minipage}[c][\hh]{\D}
      \flushleft\footnotesize Document:\\Version:\\Date:\\Page:\\ Author:
    \end{minipage}%
    \begin{minipage}[c][\hh]{\E}%
      \flushright\footnotesize\nr\\\ver\\\dat\\\pagemark\ / \pageref{LastPage}\\  \thisauthor%
    \end{minipage}%
  }
  (\textwidth, .4pt)
}
{%% Footer
  (\textwidth, 0pt)
  {} {}
  {\scriptsize \copyright Technische Universit\"at Berlin -- Chair of Space Technology\hfill}
  (\textwidth, 0pt)
}

\pagestyle{TUB}
\renewcommand{\chapterpagestyle}{TUB}
\renewcommand{\chapterheadstartvskip}{}
\setcounter{secnumdepth}{3}
\setcounter{tocdepth}{2}


% Tabellen %%%%%%%%%%%%%%%%%%%%%%%%%%%%%%%%%%%%%%%%%%%%%%%%%%%%%%%%%%%%%%%%%%%%

\usepackage{multirow}
\usepackage{tabularx}
\usepackage{longtable}   % mehrseitige Tabellen
\usepackage{booktabs}    % verschiedene Linien für Tabellen
\usepackage{threeparttable}
\newcommand{\mc}[3]{\multicolumn{#1}{#2}{#3}}

%\usepackage[%
%%parse-numbers = false,
%%parse-units = false,
%locale=DE,
%%mode=math,
%version-1-compatibility,
%group-separator = {.},
%group-digits = true,
%group-minimum-digits = 4,
%range-units=single,
%range-phrase = {--},
%retain-explicit-plus]{siunitx}
%\DeclareSIUnit\year{a}
%\DeclareSIUnit\byte{Byte}
%\DeclareSIUnit\bit{bit}
%\DeclareSIUnit\krad{\kilo\radian}
%\DeclareSIUnit\mrad{\milli\radian}
%\DeclareSIUnit\LSB{LSB}
%\DeclareSIUnit\grms{GRMS}
%\DeclareSIUnit\arcminute{arcmin}
%\DeclareSIUnit\arcsecond{arcsec}
%\DeclareSIUnit\degreeAlt{degree}
%\DeclareSIUnit\kibi{kibi}
%\DeclareSIUnit\mebi{mebi}
%\newcommand{\SIdimensions}[2]{\SIlist[list-separator={x},list-final-separator={x},list-units=single]{#1}{#2}}


% Kennzeichnung als Entwurf %%%%%%%%%%%%%%%%%%%%%%%%%%%%%%%%%%%%%%%%%%%%%%%%%%%

%\usepackage[firstpage]{draftwatermark}
%\SetWatermarkText{Draft}
%\SetWatermarkLightness{.9}


% Silbentrennung %%%%%%%%%%%%%%%%%%%%%%%%%%%%%%%%%%%%%%%%%%%%%%%%%%%%%%%%%%%%%%

\hyphenation{%
BEESAT
TUBIN
TUBSAT
wis-sen-schaft-lich
INVAR-Ka-me-ra-hal-te-run-gen
SEAR
ROS
}


% Nummerierungen und Aufzählungen %%%%%%%%%%%%%%%%%%%%%%%%%%%%%%%%%%%%%%%%%%%%%

\renewcommand{\labelenumi}{\arabic{enumi}.}
\renewcommand{\labelenumii}{\arabic{enumi}.\arabic{enumii}}
\renewcommand{\labelenumiii}{\alph{enumiii}}
\renewcommand{\labelitemi}{--}
\renewcommand{\labelitemii}{--}


% Sonderzeichen %%%%%%%%%%%%%%%%%%%%%%%%%%%%%%%%%%%%%%%%%%%%%%%%%%%%%%%%%%%%%%%

\newcommand{\grad}{\ensuremath{^\circ}} % Gradzeichen verwenden
\usepackage{eurosym}


% TODO notes %%%%%%%%%%%%%%%%%%%%%%%%%%%%%%%%%%%%%%%%%%%%%%%%%%%%%%%%%%%%%%%%%

\usepackage{todonotes}
\usepackage{xargs}

\makeatletter
\def\namedlabel#1#2{\begingroup
   \def\@currentlabel{#2}%
   \label{#1}\endgroup
}
\makeatother

\newcounter{comment}
\newcommandx*\comment[3][1=, 3=cmt:\thecomment]{
  \refstepcounter{comment}%
  {%
    \setstretch{0.7}% spacing
    \todo[inline]{\textbf{[\uppercase{#1}\thecomment]:}~#2}  \namedlabel{#3}{[#1\thecomment]}%
  }

%(#2_{#1},\ldots,#2_{#3})
}
%
%\newcommand{\comment}[2][]{%
  %% initials of the author (optional) + note in the margin
  %\refstepcounter{comment}%
  %{%
    %\setstretch{0.7}% spacing
    %\todo[inline]{%
      %\textbf{[\uppercase{#1}\thecomment]:}~#2}%
  %}
%}


% Randbedingungen %%%%%%%%%%%%%%%%%%%%%%%%%%%%%%%%%%%%%%%%%%%%%%%%%%%%%%%%%%%%%

\newlength{\G}
\newlength{\I}
\newlength{\J}
\newlength{\K}

\newcommand{\Documents}[1]{\ifthenelse{\equal{#1}{}\OR\equal{#1}{--}}{--}{#1}}

%\newcommand{\Randbedingung}[3]{%
  %\settowidth{\I}{Nummer:}
  %\settowidth{\J}{\qquad}
  %\settowidth{\K}{Dokumente:}
  %\addtolength{\I}{\J}
  %\addtolength{\I}{-\K}
  %\vspace{.75\baselineskip}\noindent\phantomsection\namedlabel{#1}{#1}%
  %\begin{tabularx}{\textwidth}{@{}X@{}}
    %\toprule
    %Nummer:\qquad\textbf{{#1}}\\
    %Dokumente:\hspace{\I}{\Dokumente{#3}}\\
    %\cmidrule{1-1}
    %{#2}\\
    %\bottomrule
  %\end{tabularx}%
  %\vspace{.75\baselineskip}}%
	
\makeatletter
\define@key{constraint}{Token}{\def\conacronym{#1}}
\define@key{constraint}{Label}{\def\conlabel{#1}}
\define@key{constraint}{Text}{\def\context{#1}}
\define@key{constraint}{Documents}{\def\condocuments{#1}}
\makeatother
\setkeys{constraint}{Label=--,Text=--,Documents=--}

\newcommand{\constraint}[1][]{%  
	\setkeys{constraint}{#1}%
	\refstepcounter{\conacronym_constraintcnt}%
  \setlength{\G}{\textwidth}%
  \settowidth{\I}{Number:}%
  \settowidth{\J}{\qquad}%
  \settowidth{\K}{Documents:}%
  \addtolength{\I}{\J}%
  \addtolength{\I}{-\K}%
  \vspace{.75\baselineskip}\noindent\phantomsection\namedlabel{\conlabel}{\mbox{\conacronym-\!\padzeroes[3]{\decimal{\conacronym_constraintcnt}}}}%
  \begin{tabularx}{\textwidth}{@{}X@{}}%
    \toprule%
    Number:\qquad\textbf{{\conacronym-\!\padzeroes[3]{\decimal{\conacronym_constraintcnt}}}}\\%
    Documents:\hspace{\I}{\Documents{\condocuments}}\\%
    \cmidrule{1-1}%
    {\context}\\%
    \bottomrule%
  \end{tabularx}%
  \vspace{.75\baselineskip}%
	}%	


% Anforderungen %%%%%%%%%%%%%%%%%%%%%%%%%%%%%%%%%%%%%%%%%%%%%%%%%%%%%%%%%%%%%%%

%\makeatletter
%\newcommand*\Herkunft[1]{%
  %\ifthenelse{\equal{#1}{}\OR\equal{#1}{--}}{--}{%
    %\@for\@tempa:=#1\do{%
      %\hyperref[\@tempa]{\@tempa}\quad}}
%}
%\makeatother

%\newcommand{\Anforderung}[4]{%
  %\setlength{\G}{\textwidth}
  %\settowidth{\I}{Nummer:}
  %\settowidth{\J}{\qquad}
  %\settowidth{\K}{Dokumente:}
  %\addtolength{\I}{\J}
  %\addtolength{\I}{-\K}
  %\vspace{.75\baselineskip}\noindent\phantomsection\namedlabel{#1}{#1}%
  %\begin{tabularx}{\textwidth}{@{}p{5.0cm}X@{}}
    %\toprule
    %Nummer:\qquad\textbf{{#1}}&Herkunft:\qquad{\Herkunft{#2}}\\
    %\multicolumn{2}{@{}p{\G}@{}}{Dokumente:\hspace{\I}{\Dokumente{#4}}}\\
    %\cmidrule{1-2}
    %\multicolumn{2}{@{}p{\G}@{}}{{#3}}\\
    %\bottomrule
  %\end{tabularx}%
  %\vspace{.75\baselineskip}}%

%\usepackage{xkeyval}
\usepackage{fmtcount}% http://ctan.org/pkg/fmtcount

\makeatletter
\newcommand*\origin[1]{%
  \ifthenelse{\equal{#1}{}\OR\equal{#1}{--}}{--}{%
    \@for\@tempa:=#1\do{%
      \ref{\@tempa}\quad}}
}
\makeatother

\makeatletter
\define@key{requirement}{Token}{\def\reqacronym{#1}}
\define@key{requirement}{Label}{\def\reqlabel{#1}}
\define@key{requirement}{Origin}{\def\reqorigin{#1}}
\define@key{requirement}{Text}{\def\reqtext{#1}}
\define@key{requirement}{Documents}{\def\reqdocuments{#1}}
\makeatother
\setkeys{requirement}{Token=REQ,Label=--,Origin=--,Text=--,Documents=--}

\newcommand{\requirement}[1][]{%  
	\setkeys{requirement}{#1}
	\refstepcounter{\reqacronym_requirementcnt}
  \setlength{\G}{\textwidth}
  \settowidth{\I}{Number:}
  \settowidth{\J}{\qquad}
  \settowidth{\K}{Documents:}
  \addtolength{\I}{\J}
  \addtolength{\I}{-\K}
  \vspace{.75\baselineskip}\noindent\phantomsection\namedlabel{\reqlabel}{\mbox{\reqacronym-\!\padzeroes[3]{\decimal{\reqacronym_requirementcnt}}}}%
  \begin{tabularx}{\textwidth}{@{}p{5.0cm}X@{}}
    \toprule
    Nummer:\qquad\textbf{{\reqacronym-\!\padzeroes[3]{\decimal{\reqacronym_requirementcnt}}}}&Herkunft:\qquad{\origin{\reqorigin}}\\
    \multicolumn{2}{@{}p{\G}@{}}{Documents:\hspace{\I}{\Documents{\reqdocuments}}}\\
    \cmidrule{1-2}
    \multicolumn{2}{@{}p{\G}@{}}{{\reqtext}}\\
    \bottomrule
  \end{tabularx}%
  \vspace{.75\baselineskip}
	}%



% hyperref.sty %%%%%%%%%%%%%%%%%%%%%%%%%%%%%%%%%%%%%%%%%%%%%%%%%%%%%%%%%%%%%%%%

\usepackage{ifpdf}
\ifpdf
\usepackage[pdfa, final]{hyperref}
%% \PassOptionsToPackage{pdfa}{hyperref}
\else
\usepackage[pdfa, dvipdfmx, final]{hyperref}
%% \PassOptionsToPackage{pdfa, dvipdfmx}{hyperref}
\fi

%% \usepackage[a-1b]{pdfx}   %% pdfx.sty lädt hyperref.sty
%% \Title{\st}
%% \Subject{Technischer Bericht}
%% \Author{TU Berlin}
%% \Keywords{TUBIN}

\hypersetup{bookmarksnumbered=true}
\hypersetup{breaklinks=true}
\hypersetup{colorlinks=true}
\hypersetup{linkcolor=blue}
\hypersetup{anchorcolor=blue}
\hypersetup{citecolor=blue}
\hypersetup{filecolor=blue}
\hypersetup{menucolor=blue}
\hypersetup{runcolor=blue}
\hypersetup{urlcolor=blue}
%% \hypersetup{pdfstartview=FitH}
\hypersetup{pdftitle={SIERRA WPD}}
\hypersetup{pdfsubject={Technical Note}}
\hypersetup{pdfauthor={TU Berlin}}
\hypersetup{pdfkeywords={SIERRA, Space, Amateur Radio, Planetary Rover}}
\hypersetup{pdflang=de-DE-1996}
\urlstyle{same}
%%\usepackage{breakurl}


% Referenzierung von Missionszielen, Randbedingungen, Anforderungen:
\newcommand{\aref}[1]{\hyperref[#1]{#1}}


% nur ein Durchlauf anstatt zwei für Package "hyperref.sty"
%\usepackage[final]{bookmark}


% level des TOC beschrŠnken
\setcounter{tocdepth}{1}


% Package "ellipsis.sty" muss nach "hyperref.sty" geladen werden (siehe README)
\usepackage[xspace]{ellipsis}


% Abkürzungsverzeichnis und Glossar %%%%%%%%%%%%%%%%%%%%%%%%%%%%%%%%%%%%%%%%%%%
% Package "glossaries.sty" muss nach "hyperref.sty" geladen werden

\usepackage[acronym, shortcuts]{glossaries}
\renewcommand*{\glspostdescription}{}  %% Punkt am Ende deaktivieren

% Zeilenabstände verkleinern
\renewcommand\glossarypreamble{\markboth{\glossaryname}{\glossaryname}}
%%\setlength{\glsdescwidth}{\linewidth}
%%\addtolength{\glsdescwidth}{-\glsdescwidth}
%%\setlength{\glspagelistwidth}{\textwidth}
\renewcommand{\glsnamefont}[1]{\textrm{#1}}
\makeglossaries

% PDF Version 5 verwenden  %%%%%%%%%%%%%%%%%%%%%%%%%%%%%%%%%%%%%%%%%%%%%%%%%%%%
% damit PDFs mit Ver. 5 eingebunden werden können

\pdfoptionpdfminorversion=5

%%%%%%%%%%%%%%%%%%%%%%%%%%%%%%%%%%%%%%%%%%%%%%%%%%%%%%%%%%%%%%%%%%%%%%%%%%%%%%%

\setcounter{tocdepth}{3}

